\chapter{Hyperparametersuchräume}
\label{chap:anhang:suchraeume}

Die Suchräume sind in zwei Tabellen zusammengefasst: Tabelle~\ref{tab:first-run-hyperparameter-spaces} enthält die Vorexperimente, Tabelle~\ref{tab:hyperparameter-spaces} die anschließende \acs{HPO} der Phase~1. Beide Tabellen umfassen die supervised Basiskonfiguration sowie \ac{LRML}, Hoffer, Label Propagation (\texttt{iscen}) und \ac{SLADE} für Cars196 (Cars), In-Shop, \ac{CUB} und Semi-iNat-2021 (Semi-iNat). Phase~1 umfasst zusätzlich In-Shop mit 10\,\% gelabelten Klassen. In Tabelle~\ref{tab:first-run-hyperparameter-spaces} bezeichnet \emph{alle} die Datensätze, auf denen ein Verfahren in den Vorexperimenten lief. Hoffer lief dort nicht auf \ac{CUB}. Jeder Verfahrensblock ergänzt die Basiswerte oder überschreibt sie für die ausdrücklich genannten Datensätze. Nicht aufgeführte verfahrensspezifische Parameter besitzen in den ausgewerteten Dateien keinen eigenen Suchraum, woraus kein bestimmter Festwert folgt.


Die Symbole entsprechen den Verlustdefinitionen im Haupttext, in Klammern steht der Parametername der Implementierung~\cite{musgrave_pytorch_losses_code_2026}. Die ArcFace-Marge ist in Grad angegeben. In den trigonometrischen Formeln gilt $m=(\pi/180)m_{\mathrm{deg}}$~\cite{musgrave_pytorch_losses_code_2026}. $B{:}n_{\mathrm{S}}$ bezeichnet die Batchgröße und die Beispiele pro Klasse. \emph{log.} bedeutet logarithmische, \emph{unif.} kontinuierliche Gleichverteilung, \emph{kat.} kategoriale Auswahl und \emph{diskr.} gleichverteilte Auswahl aus der angegebenen diskreten Menge. \emph{fest} bezeichnet einen Suchraum mit nur einem Wert. Bei Schrittweiten werden ausschließlich erreichbare Werte aufgeführt.

\section{Vorexperimente}
\label{sec:anhang:phase0}



\begingroup
\small\setstretch{1}
\setlength{\tabcolsep}{4pt}
\renewcommand{\arraystretch}{1.12}
\begin{longtable}{@{}>{\raggedright\arraybackslash\hangindent=1em\hangafter=1}p{0.30\textwidth}>{\raggedright\arraybackslash}p{0.16\textwidth}>{\raggedright\arraybackslash}p{\dimexpr0.42\textwidth-6\tabcolsep\relax}>{\raggedright\arraybackslash}p{0.12\textwidth}@{}}
\caption[Hyperparametersuchräume der Vorexperimente]{Hyperparametersuchräume der Vorexperimente: Basiskonfiguration und verfahrensspezifische Ergänzungen bzw. Überschreibungen.}\label{tab:first-run-hyperparameter-spaces}\\
\toprule
Hyperparameter & Datensatz & Suchraum & Verteilung \\
\midrule
\endfirsthead
\multicolumn{4}{l}{\tablename \thetable: Vorexperimente (Fortsetzung)}\\
\toprule
Hyperparameter & Datensatz & Suchraum & Verteilung \\
\midrule
\endhead
\midrule
\multicolumn{4}{r}{\emph{Fortsetzung auf der nächsten Seite}}\\
\endfoot
\bottomrule
\endlastfoot
\addlinespace
\multicolumn{4}{@{}l}{\emph{Supervised Basiskonfiguration}}\\*
\midrule
Klassifikator-Lernrate & alle & $[10^{-4},\,100]$ & log. \\
Lernrate & alle & $[10^{-5},\,10]$ & log. \\
Batchstruktur $B{:}n_{\mathrm{S}}$ & Cars, CUB & \{32, 64, 128, 256\}:4 & kat. \\
Batchstruktur $B{:}n_{\mathrm{S}}$ & In-Shop & \{32, 64, 128, 256, 512\}:2 & kat. \\
Batchstruktur $B{:}n_{\mathrm{S}}$ & Semi-iNat & \{128, 256, 512\}:2 & kat. \\
ArcFace $m$ [Grad] (\texttt{margin}) & alle & $[15,\,75]$ & unif. \\
ArcFace $\gamma$ (\texttt{scale}) & alle & $[10,\,250]$ & log. \\
Circle $\gamma$ (\texttt{gamma}) & alle & $[32,\,512]$ & log. \\
Circle $m$ & alle & $[0{,}01,\,0{,}5]$ & unif. \\
Multi-Similarity $\gamma^{+}$ (\texttt{alpha}) & alle & $[2,\,25]$ & unif. \\
Multi-Similarity $m$ (\texttt{base}) & alle & $[0{,}3,\,0{,}9]$ & unif. \\
Multi-Similarity $\gamma^{-}$ (\texttt{beta}) & alle & $[20,\,160]$ & log. \\
MS-Mining $m_{\mathrm{Miner}}$ (\texttt{epsilon}) & alle & $[0{,}1,\,0{,}9]$ & unif. \\
NT-Xent $\tau$ (\texttt{temperature}) & alle & $[0{,}001,\,0{,}5]$ & log. \\
Proxy-Anchor $\gamma$ (\texttt{alpha}) & alle & $[1,\,100]$ & log. \\
Proxy-Anchor $m$ (\texttt{margin}) & alle & $[0{,}01,\,1]$ & unif. \\
\addlinespace
\multicolumn{4}{@{}l}{\emph{LRML}}\\*
\midrule
Kanten-Batchgröße & Cars, In-Shop, CUB & $\{16,\,\allowbreak 32,\,\allowbreak 64,\,\allowbreak 128,\,\allowbreak 256,\,\allowbreak 512\}$ & kat. \\
Kanten-Batchgröße & Semi-iNat & $\{16,\,\allowbreak 32,\,\allowbreak 64,\,\allowbreak 128,\,\allowbreak 256,\,\allowbreak 512,\,\allowbreak 1024,\,\allowbreak 2048\}$ & kat. \\
Nachbarzahl $k_{\mathrm{NN}}$ (Regularisierer) & Cars, In-Shop, CUB & $\{3,\,\allowbreak 5,\,\allowbreak 10,\,\allowbreak 50,\,\allowbreak 100,\,\allowbreak 200\}$ & kat. \\
Nachbarzahl $k_{\mathrm{NN}}$ (Regularisierer) & Semi-iNat & $\{5,\,\allowbreak 10,\,\allowbreak 50,\,\allowbreak 100,\,\allowbreak 200,\,\allowbreak 500\}$ & kat. \\
Zielverhältnis $\eta$ & Cars, In-Shop, CUB & $[0{,}5,\,1{,}25]$ & unif. \\
Zielverhältnis $\eta$ & Semi-iNat & $[0{,}25,\,1{,}25]$ & unif. \\
Warm-up [Epochen] & Cars, In-Shop & $\{0,\,\allowbreak 5,\,\allowbreak 10,\,\allowbreak 15,\,\allowbreak 20,\,\allowbreak 25,\,\allowbreak 30\}$ & diskr. \\
Warm-up [Epochen] & CUB & $\{0,\,\allowbreak 2,\,\allowbreak 4,\,\allowbreak 6,\,\allowbreak 8,\,\allowbreak 10,\,\allowbreak 12\}$ & diskr. \\
Warm-up [Epochen] & Semi-iNat & $\{0\}$ & fest \\
\addlinespace
\multicolumn{4}{@{}l}{\emph{Hoffer}}\\*
\midrule
Zielverhältnis $\eta$ & Cars & $[0{,}5,\,1{,}25]$ & unif. \\
Zielverhältnis $\eta$ & In-Shop & $[0{,}25,\,1{,}5]$ für 36 bis 49 der 50 Versuche je Suche, $[0{,}5,\,1{,}25]$ für die übrigen & unif. \\
Zielverhältnis $\eta$ & Semi-iNat & $[0{,}5,\,1{,}5]$ & unif. \\
Warm-up [Epochen] & Cars, In-Shop & $\{0,\,\allowbreak 5,\,\allowbreak 10,\,\allowbreak 15,\,\allowbreak 20,\,\allowbreak 25,\,\allowbreak 30\}$ & diskr. \\
Distanzskala $\gamma_{\mathrm{H}}$ & alle & $[0{,}1,\,100]$ & log. \\
Referenzmengen $K_{\mathrm{ref}}$ & Cars & $\{1,\,\allowbreak 2,\,\allowbreak 3,\,\allowbreak 5\}$ & kat. \\
Referenzmengen $K_{\mathrm{ref}}$ & In-Shop, Semi-iNat & $\{1,\,\allowbreak 2,\,\allowbreak 3\}$ & kat. \\
Ungelabelte Batchgröße & alle & $\{16,\,\allowbreak 32,\,\allowbreak 64,\,\allowbreak 128\}$ & kat. \\
\addlinespace
\multicolumn{4}{@{}l}{\emph{Label Propagation (\texttt{iscen})}}\\*
\midrule
Warm-up [Epochen] & Cars, In-Shop & $\{0,\,\allowbreak 5,\,\allowbreak 10,\,\allowbreak 15,\,\allowbreak 20,\,\allowbreak 25,\,\allowbreak 30\}$ & diskr. \\
Warm-up [Epochen] & CUB & $\{0,\,\allowbreak 2,\,\allowbreak 4,\,\allowbreak 6,\,\allowbreak 8,\,\allowbreak 10,\,\allowbreak 12\}$ & diskr. \\
Konfidenzschwelle & alle & $[0{,}1,\,0{,}99]$ & unif. \\
Gelabelter Batchanteil & alle & $\{0{,}125;\,\allowbreak 0{,}25;\,\allowbreak 0{,}5\}$ & kat. \\
Diffusionsparameter $\alpha$ & alle & $[0{,}2,\,0{,}99]$ & unif. \\
Exponent $\xi$ & alle & $\{1,\,\allowbreak 2,\,\allowbreak 3,\,\allowbreak \dots,\,\allowbreak 10\}$ & diskr. \\
Nachbarzahl $k_{\mathrm{NN}}$ (Propagation) & Cars, In-Shop, CUB & $\{3,\,\allowbreak 5,\,\allowbreak 10,\,\allowbreak 50,\,\allowbreak 100,\,\allowbreak 200\}$ & kat. \\
Nachbarzahl $k_{\mathrm{NN}}$ (Propagation) & Semi-iNat & $\{5,\,\allowbreak 10,\,\allowbreak 50,\,\allowbreak 100,\,\allowbreak 200,\,\allowbreak 500\}$ & kat. \\
\addlinespace
\multicolumn{4}{@{}l}{\emph{SLADE}}\\*
\midrule
NT-Xent $\tau$ (\texttt{temperature}) & alle & $[0{,}05,\,0{,}5]$ & log. \\
Warm-up [Epochen] & Cars, In-Shop, CUB & $\{0,\,\allowbreak 5,\,\allowbreak 10,\,\allowbreak 15,\,\allowbreak 20,\,\allowbreak 25,\,\allowbreak 30\}$ & diskr. \\
Ungelabelte Batchgröße & alle & $\{32,\,\allowbreak 64,\,\allowbreak 128,\,\allowbreak 256\}$ & kat. \\
Basis: Gewicht $\beta_{\mathrm{basis}}$ des bisherigen Werts & alle & $[0{,}1,\,0{,}99]$ & unif. \\
Basis: Warm-up [Epochen] & Cars, In-Shop, CUB & $\{0;\,\allowbreak 0{,}25;\,\allowbreak 0{,}5;\,\allowbreak 1;\,\allowbreak 5\}$ & kat. \\
Basis: Warm-up [Epochen] & Semi-iNat & $\{0;\,\allowbreak 0{,}25;\,\allowbreak 0{,}5;\,\allowbreak 1;\,\allowbreak 2\}$ & kat. \\
Marge $m$ von $\mathcal{L}_{\mathrm{SD}}$ (\texttt{sd\_margin}) & alle & $[0{,}1,\,1]$ & unif. \\
Varianzgewicht $w_v$ von $\mathcal{L}_{\mathrm{SD}}$ & alle & $[0{,}5,\,20]$ & log. \\
Kontrastive Marge $m^{+}$ & alle & $[0,\,0{,}5]$ & unif. \\
Kontrastive Marge $m^{-}$ & alle & $[0{,}5,\,2]$ & unif. \\
Teacher-Aktualisierung alle $n_{\mathrm{T}}$ Epochen & Cars, In-Shop, CUB & $\{1,\,\allowbreak 10,\,\allowbreak 50,\,\allowbreak 100\}$ & kat. \\
Teacher-Aktualisierung alle $n_{\mathrm{T}}$ Epochen & Semi-iNat & $\{1,\,\allowbreak 5,\,\allowbreak 10,\,\allowbreak 20\}$ & kat. \\
\end{longtable}
\endgroup




\section{Phase 1: Hyperparameteroptimierung}
\label{sec:anhang:phase1}

Grundlage dieses Abschnitts sind die gespeicherten Konfigurationen der Phase-1-Suchen, auf denen die Ergebnisse beruhen: je Datensatz fünf Verfahren mit je fünf Verlustfunktionen. Tabelle~\ref{tab:phase1:hpo-protokoll} fasst das Protokoll der Suche zusammen und Tabelle~\ref{tab:hyperparameter-spaces} die Suchräume.

\begin{table}[htbp]
	\centering
	\caption[Protokoll der Phase-1-\acs{HPO}]{Protokoll der Phase-1-\acs{HPO} je Datensatz. In-Shop 10\,\%: 10\,\% der Trainingsklassen vollständig gelabelt, übrige Klassen ungelabelt. \UEBL: supervised Referenz. Obergrenze: Begrenzung des Gradientenverhältnisses (Abschnitt \ref{sec:m:gewicht}).}
	\label{tab:phase1:hpo-protokoll}
	\setlength{\tabcolsep}{4pt}
	\begin{tabular}{@{}>{\raggedright\arraybackslash}p{3.2cm}>{\centering\arraybackslash}p{1.9cm}>{\centering\arraybackslash}p{1.9cm}>{\centering\arraybackslash}p{1.9cm}>{\centering\arraybackslash}p{2.3cm}>{\centering\arraybackslash}p{2.0cm}@{}}
		\toprule
		Merkmal & Cars196 & CUB & In-Shop & In-Shop 10\,\% & Semi-iNat \\
		\midrule
		Label-Budget der \acs{HPO} & $k=16$ & $k=16$ & $k=2$ & 10\,\% der Klassen & $k=16$ \\
		\addlinespace
		Folds (klassendisjunkt) & 4 & 5 & 4 & 4 & 4 \\
		\addlinespace
		Versuche je Suche & \multicolumn{5}{>{\centering\arraybackslash}p{\dimexpr10.0cm+8\tabcolsep\relax}}{80: zuerst 40 zufällige Versuche, danach 40 \acs{TPE}-Versuche} \\
		\addlinespace
		Sampler & \multicolumn{5}{>{\centering\arraybackslash}p{\dimexpr10.0cm+8\tabcolsep\relax}}{Optuna-\acs{TPE}, multivariat und gruppiert} \\
		\addlinespace
		Pruning & \multicolumn{5}{>{\centering\arraybackslash}p{\dimexpr10.0cm+8\tabcolsep\relax}}{keines} \\
		\addlinespace
		Seeds & \multicolumn{5}{>{\centering\arraybackslash}p{\dimexpr10.0cm+8\tabcolsep\relax}}{7 für \acs{HPO}, Training und Auswahl der gelabelten Beispiele} \\
		\addlinespace
		Zielgröße & \multicolumn{5}{>{\centering\arraybackslash}p{\dimexpr10.0cm+8\tabcolsep\relax}}{Validierungs-\acs{MAP@R}: beste Epoche je Fold, Mittel über die Folds} \\
		\addlinespace
		Max.\ Epochen (Patience) & \multicolumn{4}{>{\centering\arraybackslash}p{\dimexpr8.0cm+6\tabcolsep\relax}}{höchstens 600 (20)} & höchstens 600 (10) \\
		\addlinespace
		Obergrenze des Gradienten\-verhältnisses & \multicolumn{5}{>{\centering\arraybackslash}p{\dimexpr10.0cm+8\tabcolsep\relax}}{\acs{LRML}: keine, Hoffer: 3, \acs{SLADE}: 3} \\
		\addlinespace
		Gesuchte Hyperparameter je Suche & \multicolumn{5}{>{\centering\arraybackslash}p{\dimexpr10.0cm+8\tabcolsep\relax}}{\UEBL: 3--6, \acs{LP}: 9--12, \acs{LRML}: 7--10, Hoffer: 8--11, \acs{SLADE}: 13--17} \\
		\bottomrule
	\end{tabular}
\end{table}
 %


So ist Tabelle~\ref{tab:hyperparameter-spaces} zu lesen: Ist der Suchraum auf allen Datensätzen gleich, steht er zentriert über allen Spalten. Getrennte Zellen zeigen, wo sich die Datensätze unterscheiden. Ein Intervall ohne Zusatz ist kontinuierlich gleichverteilt, \emph{log.} kennzeichnet eine logarithmische Gleichverteilung. Mengen werden gleichverteilt ausgewählt (ganzzahlige Parameter mit Schrittweite und kategoriale Parameter). \emph{fest} bezeichnet eine Menge mit nur einem Wert. $B{:}n_{\mathrm{S}}$ bezeichnet die Batchgröße und die Beispiele pro Klasse. Bei ArcFace und Proxy-Anchor werden die Lernrate und die Klassifikator-Lernrate unabhängig voneinander gesucht. \emph{wie oben} verweist auf die Zeile im Block der Verlustfunktionen.

\begingroup
\newcommand*{\spVier}[1]{\multicolumn{4}{>{\centering\arraybackslash}p{\dimexpr0.62\textwidth-20.8pt+6\tabcolsep\relax}}{#1}}
\newcommand*{\spFuenf}[1]{\multicolumn{5}{>{\centering\arraybackslash}p{\dimexpr0.775\textwidth-26pt+8\tabcolsep\relax}}{#1}}
\small\setstretch{1}
\setlength{\tabcolsep}{2.5pt}
\renewcommand{\arraystretch}{1.15}
\begin{longtable}{@{}>{\raggedright\arraybackslash\hangindent=1em\hangafter=1}p{0.22\textwidth}*{5}{>{\centering\arraybackslash}p{\dimexpr0.155\textwidth-5.2pt\relax}}@{}}
\caption[Hyperparametersuchräume der Phase 1]{Hyperparametersuchräume der Phase 1, gelesen aus den Konfigurationen der ausgewerteten Suchen.}\label{tab:hyperparameter-spaces}\\
\toprule
Hyperparameter & Cars196 & CUB & In-Shop & In-Shop 10\,\% & Semi-iNat \\
\midrule
\endfirsthead
\multicolumn{6}{@{}l}{\small Tabelle \thetable{} (Fortsetzung)}\\
\toprule
Hyperparameter & Cars196 & CUB & In-Shop & In-Shop 10\,\% & Semi-iNat \\
\midrule
\endhead
\midrule
\multicolumn{6}{r@{}}{\emph{Fortsetzung auf der nächsten Seite}}\\
\endfoot
\bottomrule
\endlastfoot
\addlinespace
\multicolumn{6}{@{}l}{\emph{Alle Verfahren außer SLADE: Optimierung und Batch}}\\*
\midrule
Batchstruktur $B{:}n_{\mathrm{S}}$ & \{32, 64, 128, 256\}:4 & \{32, 64, 128, 256\}:4 & \{64, 128, 256, 512\}:2 & \{64, 128, 256, 512\}:8 & \{128, 256, 512\}:2 \\
Lernrate (Circle, MS, NT-Xent) & $[10^{-5},\allowbreak\,5\cdot 10^{-4}]$ log. & $[5\cdot 10^{-6},\allowbreak\,10^{-4}]$ log. & $[10^{-5},\allowbreak\,5\cdot 10^{-4}]$ log. & $[10^{-5},\allowbreak\,5\cdot 10^{-4}]$ log. & $[5\cdot 10^{-6},\allowbreak\,10^{-4}]$ log. \\
Lernrate (ArcFace, Proxy-Anchor) & \spVier{$[10^{-5},\,10]$ log.} & $[5\cdot 10^{-6},\allowbreak\,10^{-4}]$ log. \\
Klassifikator-Lernrate (ArcFace, Proxy-Anchor) & \spVier{$[10^{-5},\,10]$ log.} & $[10^{-5},\,100]$ log. \\
\addlinespace
\multicolumn{6}{@{}l}{\emph{Verlustfunktionen (alle Verfahren, Abweichungen bei SLADE siehe unten)}}\\*
\midrule
ArcFace $\gamma$ (\texttt{scale}) & \spFuenf{$[0{,}1,\,80]$ log.} \\
ArcFace $m$ [Grad] (\texttt{margin}) & \spVier{$[1,\,80]$} & $[1,\,70]$ \\
Circle $\gamma$ (\texttt{gamma}) & \spFuenf{$[1,\,512]$ log.} \\
Circle $m$ & \spFuenf{$[0{,}01,\,0{,}5]$ log.} \\
Multi-Similarity $\gamma^{+}$ (\texttt{alpha}) & \spFuenf{$[0{,}1,\,30]$ log.} \\
Multi-Similarity $\gamma^{-}$ (\texttt{beta}) & \spFuenf{$[1,\,280]$ log.} \\
Multi-Similarity $m$ (\texttt{base}) & \spVier{$[0{,}4,\,0{,}99]$} & $[0{,}01,\,0{,}7]$ \\
MS-Mining $m_{\mathrm{Miner}}$ (\texttt{epsilon}) & \spFuenf{$[0{,}1,\,0{,}99]$} \\
NT-Xent $\tau$ (\texttt{temperature}) & \spVier{$[0{,}01,\,0{,}4]$ log.} & $[0{,}005,\,0{,}15]$ log. \\
Proxy-Anchor $\gamma$ (\texttt{alpha}) & \spVier{$[0{,}1,\,30]$ log.} & $[1,\,80]$ log. \\
Proxy-Anchor $m$ (\texttt{margin}) & \spVier{$[0{,}01,\,1]$ log.} & $[0{,}01,\,1]$ \\
\addlinespace
\multicolumn{6}{@{}l}{\emph{Label Propagation (\texttt{iscen})}}\\*
\midrule
Warm-up [Epochen] & \{0, 5, \dots, 30\} & \{0, 2, \dots, 12\} & \{0, 5, \dots, 30\} & \{0, 5, \dots, 30\} & \{0, 2, \dots, 10\} \\
Nachbarzahl $k_{\mathrm{NN}}$ & \{2, 3, \dots, 100\} & \{2, 3, \dots, 200\} & \{2, 3, \dots, 100\} & \{2, 3, \dots, 100\} & \{2, 3, \dots, 100\} \\
Diffusions\-parameter $\alpha$ & \spFuenf{$[0{,}1,\,0{,}99]$} \\
Exponent $\xi$ & \spFuenf{\{1, 2, \dots, 10\}} \\
Konfidenzschwelle & \spFuenf{$[0{,}1,\,0{,}9]$} \\
Gelabelter Batchanteil & \spFuenf{\{0,125; 0,25; 0,5\}} \\
\addlinespace
\multicolumn{6}{@{}l}{\emph{LRML}}\\*
\midrule
Warm-up [Epochen] & \{0, 5, \dots, 30\} & \{0, 2, \dots, 12\} & \{0, 5, \dots, 30\} & \{0, 5, \dots, 30\} & \{0, 1, 2\} \\
Zielverhältnis $\eta$ & \spVier{$[0{,}5,\,1{,}5]$} & $[0{,}5,\,1{,}25]$ \\
Nachbarzahl $k_{\mathrm{NN}}$ & \{2, 3, \dots, 100\} & \{2, 3, \dots, 200\} & \{2, 3, \dots, 50\} & \{2, 3, \dots, 50\} & \{2, 3, \dots, 100\} \\
Kanten-Batchgröße & \spVier{\{16, 32, 64, 128, 256, 512\}} & \{16, 32, 64\} \\
\addlinespace
\multicolumn{6}{@{}l}{\emph{Hoffer}}\\*
\midrule
Warm-up [Epochen] & \{0, 5, \dots, 30\} & \{0, 2, \dots, 12\} & \{0, 5, \dots, 30\} & \{0, 5, \dots, 30\} & \{0, 1, 2\} \\
Zielverhältnis $\eta$ & \spVier{$[0{,}5,\,1{,}5]$} & $[0{,}5,\,1{,}25]$ \\
Distanzskala $\gamma_{\mathrm{H}}$ & $[1,\,30]$ log. & $[1,\,30]$ log. & $[1,\,30]$ log. & $[1,\,100]$ & $[4,\,256]$ log. \\
Referenzmengen $K_{\mathrm{ref}}$ & \{1, 2, 3, 5\} & \{1, 2, 3, 5\} & \{1\} (fest) & \{1\} (fest) & \{1\} (fest) \\
Ungelabelte Batchgröße & \{16, 32, 64, 128\} & \{16, 32, 64, 128\} & \{16, 32, 64\} & \{16, 32, 64\} & \{16, 32, 64, 128\} \\
\addlinespace
\multicolumn{6}{@{}l}{\emph{SLADE}}\\*
\midrule
Batchstruktur $B{:}n_{\mathrm{S}}$ & \{32, 64, 128, 256\}:4 & \{64, 128, 256\}:4 & \{128, 256, 512\}:2 & \{128, 256, 512\}:8 & \{128, 256, 512\}:2 \\
Lernrate (alle Verluste) & $[10^{-5},\allowbreak\,5\cdot 10^{-4}]$ log. & $[5\cdot 10^{-6},\allowbreak\,10^{-4}]$ log. & $[10^{-5},\allowbreak\,5\cdot 10^{-4}]$ log. & $[10^{-5},\allowbreak\,5\cdot 10^{-4}]$ log. & $[5\cdot 10^{-6},\allowbreak\,10^{-4}]$ log. \\
Klassifikator-Lernrate (ArcFace, Proxy-Anchor) & \spFuenf{$[10^{-5},\,10]$ log.} \\
Multi-Similarity $m$ (\texttt{base}) (abweichend) & \spFuenf{$[0{,}4,\,0{,}9]$} \\
NT-Xent $\tau$ (\texttt{temperature}) (abweichend) & \spVier{$[0{,}05,\,0{,}3]$ log.} & $[0{,}02,\,0{,}15]$ log. \\
ArcFace $m$ [Grad] (\texttt{margin}) (abweichend) & \spVier{wie oben} & $[1,\,80]$ \\
Proxy-Anchor $\gamma$ (\texttt{alpha}) (abweichend) & \spVier{wie oben} & $[0{,}1,\,30]$ log. \\
Proxy-Anchor $m$ (\texttt{margin}) (abweichend) & \spVier{wie oben} & $[0{,}01,\,1]$ log. \\
Warm-up [Epochen] & \{0, 5, \dots, 30\} & \{0, 2, \dots, 12\} & \{0, 5, \dots, 30\} & \{0, 5, \dots, 30\} & \{0, 1, 2\} \\
Teacher-Aktualisierung alle $n_{\mathrm{T}}$ Epochen & \{1, 5, 10\} & \{1, 2, 5\} & \{1, 5, 10\} & \{1, 5, 10\} & \{1, 2, 5\} \\
Zielverhältnis $\eta$ & \spVier{$[0{,}5,\,1{,}5]$} & $[0{,}5,\,1{,}25]$ \\
Clusterverhältnis (Cluster je gelabelter Klasse) & \{1, 2, 3\} & \{1, 2, 3\} & \{0,1; 0,5; 1; 2\} & \{0,1; 0,5; 1; 2\} & \{1, 2, 3\} \\
Ungelabelte Batchgröße & \spVier{\{32, 64, 128, 256, 512\}} & \{32, 64, 128\} \\
Basis: Zielverhältnis $\eta_{\mathrm{basis}}$ & \spVier{$[0{,}5,\,1{,}5]$} & $[0{,}5,\,1{,}25]$ \\
Basis: Warm-up [Epochen] & \{0; 0,25; 0,5; 1; 5\} & \{0; 0,25; 0,5; 1; 2,5\} & \{0; 0,25; 0,5; 1; 5\} & \{0; 0,25; 0,5; 1; 5\} & \{0; 0,25; 0,5; 1\} \\
Basis: Gewicht $\beta_{\mathrm{basis}}$ des bisherigen Werts & \spFuenf{$[0{,}1,\,0{,}99]$} \\
$\mathcal{L}_{\mathrm{SD}}$: Marge $m$ (\texttt{sd\_margin}) & \spFuenf{$[0{,}1,\,0{,}99]$} \\
$\mathcal{L}_{\mathrm{SD}}$: Varianzgewicht $w_v$ & \spFuenf{$[0{,}1,\,20]$ log.} \\
Kontrastive Marge $m^{+}$ (ArcFace, Proxy-Anchor) & \spFuenf{$[0,\,0{,}5]$} \\
Kontrastive Marge $m^{-}$ (ArcFace, Proxy-Anchor) & \spFuenf{$[0{,}5,\,2]$} \\
\end{longtable}
\endgroup






